\scanpage{112}
\mspage{61}
$f_{U'},f_{T'}$ :
\diagram{\begin{tikzcd}[ampersand replacement=\&]
U'\arrow[r,hook,no head,"i'"]\arrow[d,"{f_{U'}}"'] \& X'\arrow[d,"f"] \& T'\arrow[l,hook',no head,"j'"']\arrow[d,"{f_{T'}}"] \\
U\arrow[r,hook,"i"] \& X \& T\arrow[l,hook',"j"']\quad ;
\end{tikzcd}}
elle s'obtient à l'aide du bifoncteur
\[
\left\{\begin{aligned} H'(u',t')&\overset{\mathrm{def}}{=}H(f_{U'}(u'),f_{T'}(t'))\\ \text{i.e.}\qquad H'&= H\circ(f_{U'}^{\unsure{\circ|o}}\times f_{T'}).\end{aligned}\right.
\]
\unsure{Cet|Ce} triple \inserted{i.e. cet objet de $\unsure{\mathscr T|\mathcal T}$} représente le foncteur suivant sur $\mathscr T$ :
\[
(X'',U'',T'')\longmapsto\mathrm{Hom}_{\mathrm{Cat}}(U'',U)\times\mathrm{Hom}_{\mathrm{Cat}}(T'',T),
\]
i.e. on a pour tout objet $(X'',U'',T'')$ de $\mathscr T$ une bijection can., fonctorielle en cet objet
\[
\mathrm{Hom}_{\mathscr T}((X'',U'',T''),(X',U',T'))\simeq\mathrm{Hom}_{\mathrm{Cat}}(U'',U)\times\mathrm{Hom}_{\mathrm{Cat}}(T'',T).\tag*{(*)}
\]
Nul doute qu'on a même
\[
\underline{\mathrm{Hom}}_{\mathscr T}((X'',U'',T''),(X',U',T'))\simeq\underline{\mathrm{Hom}}(U'',\unsure{U|U'})\times\underline{\mathrm{Hom}}(T'',\unsure{T|T'}),\tag*{(**)}
\]
tout aussi fonctoriel, dont (*) se déduit en passant aux ens. d'objets. On dit que le triple $(X',U',T'')$ \deleted{$X'$ est \unsure{dé}} est \deleted{déduit} le triple \uline{induit} par le triple $(X,U,T)$, via les foncteurs $f_{\unsure{U'|U}},f_{T'}$.
\margin{\uline{NB} Si $f_{U'},f_{T'}$ sont pl. fidèles, alors il en est de même de $f:X'\longrightarrow X$. Nous utiliserons cette opération de recollement privilégié surtout quand $U'=U,f_{U'}=\mathrm{id}_U$, ou $T'=T,f_{T'}=\mathrm{id}_T$.}

\deleted{\unsure{Rel}} 1.2. \uline{Relation avec le recollement des topos.}
La situation d'un triple \deleted{d'un}
\[
(\unsure{\mathscr X|\mathcal X},\mathscr U,\mathscr T)
\]
d'un topos, $\mathscr X$, d'un ouvert $U$ de $\mathscr X$, et d'une
% archive sheet number: 112
